\section{最小二乘与数学结构}\label{sec:math}
设 \(\vect{A}\in\R^{m\times n}\)、\(\vect{x}\in\R^n\)、
\(\vect{b}\in\R^m\)，所有元素均为实数。最小二乘目标为
\begin{equation}
  J(\vect{x})=\frac{1}{2}\lVert\vect{A}\vect{x}-\vect{b}\rVert_2^2.
  \label{eq:least-squares}
\end{equation}
两条分别编号的公式使用 \texttt{align}，以等号为对齐点：
\begin{align}
  \nabla J(\vect{x})&=\vect{A}^{\mathsf{T}}
    (\vect{A}\vect{x}-\vect{b}),\label{eq:gradient}\\
  \nabla^2J(\vect{x})&=\vect{A}^{\mathsf{T}}\vect{A}.
    \label{eq:hessian}
\end{align}
一条推导需要多行、但只保留一个编号时，使用外层公式包住 \texttt{aligned}：
\begin{equation}
  \begin{aligned}
    J(\vect{x})
      &=\frac{1}{2}(\vect{A}\vect{x}-\vect{b})^{\mathsf{T}}
        (\vect{A}\vect{x}-\vect{b})\\
      &=\frac{1}{2}\vect{x}^{\mathsf{T}}\vect{A}^{\mathsf{T}}
        \vect{A}\vect{x}
        -\vect{b}^{\mathsf{T}}\vect{A}\vect{x}
        +\frac{1}{2}\vect{b}^{\mathsf{T}}\vect{b}.
  \end{aligned}
  \label{eq:expanded-loss}
\end{equation}
以阈值 \(\delta>0\) 定义分段损失，条件列与表达式列分开对齐：
\begin{equation}
  \rho_\delta(r)=
  \begin{cases}
    \dfrac{1}{2}r^2,&\lvert r\rvert\leq\delta,\\
    \delta\bigl(\lvert r\rvert-\dfrac{1}{2}\delta\bigr),
      &\lvert r\rvert>\delta.
  \end{cases}
  \label{eq:huber}
\end{equation}
矩阵环境中的分列、换行同样表达结构，例如
\[
  \vect{A}=\begin{bmatrix}1&0\\1&1\\1&2\end{bmatrix},
  \qquad
  \vect{x}=\begin{bmatrix}a\\b\end{bmatrix},
  \qquad
  \vect{b}=\begin{bmatrix}1\\2\\2\end{bmatrix}.
\]
\begin{theorem}[正定二次型]\label{thm:positive-quadratic}
若实对称矩阵 \(\vect{Q}\) 正定，则
\(f(\vect{x})=\vect{x}^{\mathsf{T}}\vect{Q}\vect{x}\)
的唯一全局最小点为 \(\vect{x}=\vect{0}\)。
\end{theorem}
\begin{proof}
由正定性，非零向量满足 \(f(\vect{x})>0\)，而 \(f(\vect{0})=0\)。
\end{proof}
练习：引用\cref{eq:gradient,eq:hessian}，再插入新公式观察编号变化。
将\cref{thm:positive-quadratic}的“正定”改成“半正定”后，数学结论需要重新检查。
